\subsection{Vector Functions}

\content{
\begin{definition} A vector function is a function which takes in each element
in it's domain, and returns a vector.
\end{definition}

In general, a vector function whose domain is the real numbers, and whose range
is the set of 3-dimensional vectors takes the form 
$$ \vec{r}(t) = \la f(t), g(t), h(t)\ra. $$

\begin{example} If 
$$ \vec{r}(t) = \la t^3, \ln(3-t), \sqrt{t}\ra, $$
find the domain of $\vec{r}$. 
\end{example}
}{% presentation
\vspace{2in}
}{% nopresentation
}{% comments
By convention, the domain of a vector valued function is the set of all real
vectors for which it is defined. 
}

\content{
\begin{theorem} If $\vec{r}(t)=\la f(t),g(t),h(t)\ra$ is a vector valued
function, then
$$ \lim_{t\rightarrow a} \vec{r}(t) = \la \lim_{t\rightarrow a}f(t), 
\lim_{t\rightarrow a}g(t), \lim_{t\rightarrow a}h(t)\ra. $$
\end{theorem}

\begin{example} Find $\lim_{t\rightarrow 0} \vec{r}(t)$ where 
$$ \vec{r}(t) = \la 1+t^3, te^{-t}, \frac{\sin t}{t}\ra. $$
\end{example}
}{% presentation
\vspace{2in}
}{% nopresentation
}{% comments
}

\content{
\begin{definition} A vector function $\vec{r}(t)$ is said to be continuous at
$a$ if 
$$ \lim_{t\rightarrow a}\vec{r}(t) = \vec{r}(a). $$
\end{definition}
}{% presentation
}{% nopresentation
}{% comments
}

\content{
\begin{example} Describe the curve defined by the equation 
$$ \vec{r}(t) = \la 1+2t, 2-t, -1+6t\ra. $$
\end{example}
}{% presentation
\vspace{2in}
}{% nopresentation
}{% comments
}

\content{
\begin{example} Sketch the curve defined by 
$$\vec{r}(t) =  \la \cos t, \sin t, t\ra. $$
\end{example}
}{% presentation
\vspace{2in}
}{% nopresentation
}{% comments
}

\content{
\begin{example} Find a vector function that represents the curve of intersection
between the cylinder $x^2+y^2=1$ and the plane $y+z=2$. 
\end{example}
}{% presentation
\vspace{3in}
}{% nopresentation
}{% comments
}

\subsection{Calculus with Vector Functions}

\content{
\begin{definition} If $\vec{r}(t)$ is a vector function, it's derivative is
defined by 
$$ \vec{r}'(t) = \lim_{h\rightarrow 0} \frac{\vec{r}(t+h)-\vec{r}(t)}{h}. $$
Provided that this limit exists. 
\end{definition}
}{% presentation
}{% nopresentation
}{% comments
}

\content{
\begin{theorem} If $\vec{r}(t) = \la f(t), g(t), h(t)\ra$, and $f$, $g$, and $h$
are differentiable, then 
$$\vec{r}'(t) = \la f'(t), g'(t), h'(t)\ra. $$
\end{theorem}

\begin{definition} The unit vector with the same direction as the tangent vector
to $\vec{r}(t)$ is called the unit tangent vector and is given by 
$$ \vec{T}(t) = \frac{\vec{r}(t)}{|\vec{r}(t)|}. $$
\end{definition}
}{% presentation
}{% nopresentation
}{% comments
}

\content{
\begin{example} Find the derivative of $\vec{r}(t) = \la 1+t^3, te^{-t}, \sin
2t\ra,$ and find the unit tangent vector at the point $t=0$.
\end{example}
}{% presentation
\vspace{2in}
}{% nopresentation
}{% comments
}

\content{
\begin{theorem} Suppose $\vec{u}$ and $\vec{v}$ are differentiable vector valued
functions, and $c$ is a scalar, and $f(t)$ is a real valued function, then 
\begin{itemize}
\item $\frac{d}{dt}(\vec{u}+\vec{v}) = \vec{u}'(t) + \vec{v}'(t)$
\item $\frac{d}{dt}(c\vec{u}(t)) = c\vec{u}'(t)$
\item $\frac{d}{dt}(f(t)\vec{u}(t)) = f(t)\vec{u}'(t) + f'(t)\vec{u}(t)$
\item $\frac{d}{dt}(\vec{u}(t)\cdot \vec{v}(t)) = \vec{u}'(t)\cdot \vec{v}(t) +
\vec{u}(t)\cdot \vec{v}'(t)$
\item $\frac{d}{dt}(\vec{u}(t)\times\vec{v}(t)) = \vec{u}'(t)\times\vec{v}(t) +
\vec{u}(t)\times\vec{v}'(t)$
\item $\frac{d}{dt}(\vec{u}(f(t))) = f'(t)\vec{u}'(f(t))$
\end{itemize}
\end{theorem}
}{% presentation
}{% nopresentation
}{% comments
}

\content{
\begin{theorem} If $|\vec{r}(t)|=c$ is a constant for all $t$, then
$\vec{r}'(t)$ is orthogonal to $\vec{r}(t)$ for all $t$. 
\end{theorem}
}{% presentation
\vspace{2in}
}{% nopresentation
}{% comments
Start with $\vec{r}(t)\cdot\vec{r}(t) = c^2$, then use the previous theorem to
compute $\frac{d}{dt}(\vec{r}(t)\cdot \vec{r}(t))$.
}

\content{
\begin{example} Compute the derivative of the vector valued function defined by 
$$ \vec{r}(t) = \sin(t)\cdot \la 1,t, e^t\ra. $$
\end{example}
}{% presentation
\vspace{2in}
}{% nopresentation
}{% comments
}

\content{
The integral of a vector valued function is similarly computed: If
$\vec{r}(t)=\la f(t), g(t), h(t)\ra$, and each of $f$, $g$, and $h$ is
integrable on $[a,b]$, then 
$$ \int_a^b \vec{r}(t) dt = \la \int_a^b f(t)dt, \int_a^b g(t)dt, \int_a^b
h(t)dt \ra. $$
}{% presentation
}{% nopresentation
}{% comments
}

